% Long \code{} run names are unbreakable; \sloppy keeps them from producing
% overfull boxes. Scoped to this section only.
\begingroup
\sloppy

\section{Cross-method comparison}
\label{sec:cmp}

% AUTO-GENERATED by experiments_report/scripts/extract_comparison.py
% Do not edit by hand; re-run the script instead.

\newcommand{\CmpMasterUDALES}{%
\footnotesize
\begin{adjustbox}{max width=\textwidth}
\begin{tabular}{@{}ll rrr rr rrrr r@{}}
\toprule
ID & method & \multicolumn{3}{c}{sensor / field RMSE} & \multicolumn{2}{c}{parameters} & \multicolumn{4}{c}{calibration} & wall \\
\cmidrule(lr){3-5}\cmidrule(lr){6-7}\cmidrule(lr){8-11}
 & & assim & valid & field $\Umag$ & $\alpha$ [deg] & $\Umag$ [m/s] & $\chi^2$ & $z_{\mathrm{par}}$ & $z_{\mathrm{val}}$ & $q$ & [h] \\
\midrule
\multicolumn{12}{@{}l}{\emph{\code{inflow} (laminar), correlation loc.}} \\
E11 & ESMDA & 0.665 & 0.748 & 0.477 & 7.74 & \textbf{0.529} & -- & 12.07 & 6.16 & \textbf{0.719} & 2.19 \\
F8 & EnKF state$+$RTPS & 0.167 & 0.626 & 0.261 & \textit{19.29}$^{\dagger}$ & \textit{0.919}$^{\dagger}$ & \textbf{0.92} & \textit{1.88}$^{\dagger}$ & \textbf{1.16} & 0.478 & \textbf{1.26} \\
F7 & EnKF joint$+$RTPS & \textbf{0.155} & \textbf{0.394} & \textbf{0.247} & \textbf{5.26} & 0.910 & 0.90 & \textbf{1.76} & 0.83 & 0.606 & 1.26 \\
H11 & Filter smoothing & 0.506 & 0.593 & 0.486 & 5.81 & 0.616 & 13.26 & 10.59 & 8.75 & 0.433 & 2.87 \\
\midrule
\multicolumn{12}{@{}l}{\emph{\code{inflow} (laminar), no loc.}} \\
E18 & ESMDA & 0.694 & 0.754 & 0.521 & 7.59 & 0.671 & -- & 13.63 & 7.82 & \textbf{0.724} & 2.15 \\
F14 & EnKF state, no infl. & 0.236 & 0.570 & 0.299 & \textit{19.29}$^{\dagger}$ & \textit{0.919}$^{\dagger}$ & 3.98 & \textit{1.88}$^{\dagger}$ & 2.20 & 0.533 & 1.19 \\
F15 & EnKF state$+$RTPS & 0.157 & 0.539 & 0.229 & \textit{19.29}$^{\dagger}$ & \textit{0.919}$^{\dagger}$ & 0.95 & \textit{1.88}$^{\dagger}$ & \textbf{1.07} & 0.538 & \textbf{1.18} \\
F13 & EnKF joint$+$RTPS & \textbf{0.144} & \textbf{0.319} & \textbf{0.221} & \textbf{3.90} & 0.818 & \textbf{0.99} & \textbf{2.23} & 0.88 & 0.626 & 1.19 \\
H18 & Filter smoothing & 0.524 & 0.620 & 0.438 & 7.87 & \textbf{0.506} & 14.96 & 15.67 & 11.25 & 0.433 & 2.81 \\
\midrule
\multicolumn{12}{@{}l}{\emph{\code{inflow\_turb}, correlation loc.}} \\
E12 & ESMDA & 0.401 & \textbf{0.365} & \textbf{0.372} & \textbf{2.31} & \textbf{0.369} & -- & 4.61 & 2.71 & \textbf{0.620} & 2.30 \\
F10 & EnKF state$+$RTPS & 0.240 & 0.753 & 0.445 & \textit{19.29}$^{\dagger}$ & \textit{0.919}$^{\dagger}$ & 1.13 & \textit{1.88}$^{\dagger}$ & 1.30 & 0.442 & \textbf{1.26} \\
F9 & EnKF joint$+$RTPS & \textbf{0.206} & 0.559 & 0.422 & 3.36 & 0.912 & \textbf{1.04} & \textbf{0.74} & \textbf{0.99} & 0.485 & 1.27 \\
H12 & Filter smoothing & 0.272 & 0.497 & 0.533 & 2.61 & 0.433 & 3.68 & 5.01 & 2.36 & 0.535 & 2.93 \\
\midrule
\multicolumn{12}{@{}l}{\emph{\code{inflow\_turb}, no loc.}} \\
E19 & ESMDA & 0.456 & \textbf{0.407} & \textbf{0.421} & \textbf{2.45} & 0.439 & -- & 7.81 & 3.35 & 0.619 & 2.29 \\
F17 & EnKF state, no infl. & 0.390 & 0.991 & 0.687 & \textit{19.29}$^{\dagger}$ & \textit{0.919}$^{\dagger}$ & 8.79 & \textit{1.88}$^{\dagger}$ & 3.49 & 0.458 & 1.20 \\
F18 & EnKF state$+$RTPS & 0.261 & 0.772 & 0.476 & \textit{19.29}$^{\dagger}$ & \textit{0.919}$^{\dagger}$ & \textbf{1.82} & \textit{1.88}$^{\dagger}$ & \textbf{1.42} & 0.486 & \textbf{1.20} \\
F16 & EnKF joint$+$RTPS & 0.260 & 0.624 & 0.500 & 4.88 & 0.955 & 1.95 & \textbf{3.29} & 1.48 & 0.504 & 1.21 \\
H19 & Filter smoothing & \textbf{0.256} & 0.441 & 0.430 & 3.45 & \textbf{0.365} & 3.43 & 9.06 & 2.79 & \textbf{0.728} & 2.89 \\
\midrule
\multicolumn{12}{@{}l}{\emph{\code{periodic}, correlation loc.}} \\
E13 & ESMDA & 1.091 & 1.408 & 1.040 & \textbf{13.62} & \textbf{0.717} & -- & 2.58 & 1.47 & 0.705 & 1.77 \\
F12 & EnKF state$+$RTPS & 0.315 & 1.261 & 1.029 & \textit{19.29}$^{\dagger}$ & \textit{0.919}$^{\dagger}$ & \textbf{1.76} & \textit{1.88}$^{\dagger}$ & \textbf{1.22} & 0.832 & \textbf{1.01} \\
F11 & EnKF joint$+$RTPS & 0.317 & 1.551 & 1.733 & 19.03 & 4.221 & 1.85 & \textbf{2.22} & 1.35 & \textbf{0.887} & 1.06 \\
H13 & Filter smoothing & \textbf{0.314} & \textbf{1.246} & \textbf{1.019} & 14.58 & 0.756 & 1.81 & 2.67 & 1.40 & 0.867 & 2.38 \\
\midrule
\multicolumn{12}{@{}l}{\emph{\code{periodic}, no loc.}} \\
E20 & ESMDA & 1.163 & \textbf{1.671} & \textbf{1.075} & 14.84 & \textbf{0.764} & -- & 89.85 & \textbf{1.32} & 0.578 & 1.80 \\
F20 & EnKF state, no infl. & 0.735 & 2.391 & 1.968 & \textit{19.29}$^{\dagger}$ & \textit{0.919}$^{\dagger}$ & 27.72 & \textit{1.88}$^{\dagger}$ & 7.00 & 0.659 & \textbf{1.11} \\
F21 & EnKF state$+$RTPS & 0.707 & 2.362 & 2.061 & \textit{19.29}$^{\dagger}$ & \textit{0.919}$^{\dagger}$ & \textbf{14.05} & \textit{1.88}$^{\dagger}$ & 3.98 & 0.655 & 1.12 \\
F19 & EnKF joint$+$RTPS & 0.826 & 3.653 & 2.951 & 19.12 & 2.173 & 17.82 & \textbf{11.54} & 4.09 & 0.483 & 1.16 \\
H20 & Filter smoothing & \textbf{0.706} & 1.766 & 1.529 & \textbf{8.89} & 0.928 & 23.06 & 41.51 & 6.50 & \textbf{0.761} & 2.35 \\
\bottomrule
\end{tabular}
\end{adjustbox}
}

\newcommand{\CmpMasterPALM}{%
\footnotesize
\begin{adjustbox}{max width=\textwidth}
\begin{tabular}{@{}ll rrr rr rrrr r@{}}
\toprule
ID & method & \multicolumn{3}{c}{sensor / field RMSE} & \multicolumn{2}{c}{parameters} & \multicolumn{4}{c}{calibration} & wall \\
\cmidrule(lr){3-5}\cmidrule(lr){6-7}\cmidrule(lr){8-11}
 & & assim & valid & field $\Umag$ & $\alpha$ [deg] & $\Umag$ [m/s] & $\chi^2$ & $z_{\mathrm{par}}$ & $z_{\mathrm{val}}$ & $q$ & [h] \\
\midrule
\multicolumn{12}{@{}l}{\code{inflow} (laminar), correlation loc. --- \bad{all methods failed (PALM truth diverged)}} \\
\addlinespace[2pt]
\multicolumn{12}{@{}l}{\code{inflow} (laminar), no loc. --- \bad{all methods failed (PALM truth diverged)}} \\
\midrule
\multicolumn{12}{@{}l}{\emph{\code{inflow\_turb}, correlation loc.}} \\
E2 & ESMDA & 1.422 & 1.377 & 1.082 & 7.12 & 0.472 & -- & 12.49 & 10.73 & -- & 2.36 \\
F2 & EnKF joint$+$RTPS & \textbf{0.627} & \textbf{1.297} & 1.151 & 11.11 & 0.935 & \textbf{14.07} & \textbf{5.38} & \textbf{3.03} & -- & \textbf{1.28} \\
H2 & Filter smoothing & 0.812 & 1.429 & \textbf{0.991} & \textbf{6.96} & \textbf{0.435} & 34.35 & 12.44 & 8.64 & -- & 2.93 \\
\midrule
\multicolumn{12}{@{}l}{\emph{\code{inflow\_turb}, no loc.}} \\
E9 & ESMDA & 1.481 & \textbf{1.389} & \textbf{1.050} & \textbf{6.62} & 0.561 & -- & 18.38 & 10.14 & -- & 2.35 \\
F5 & EnKF joint$+$RTPS & \textbf{0.880} & 1.891 & 1.582 & 13.37 & 1.261 & \textbf{27.43} & \textbf{8.31} & \textbf{4.46} & -- & \textbf{1.21} \\
H9 & Filter smoothing & 0.942 & 1.416 & 1.280 & 6.89 & \textbf{0.458} & 43.27 & 15.43 & 10.67 & -- & 2.96 \\
\midrule
\multicolumn{12}{@{}l}{\emph{\code{periodic}, correlation loc.}} \\
E3 & ESMDA & 1.437 & 2.133 & \textbf{4.014} & 22.42 & \textbf{1.090} & -- & 3.13 & 2.77 & -- & 2.08 \\
F3 & EnKF joint$+$RTPS & 0.704 & 4.072 & 5.905 & 138.50 & 21.276 & 15.86 & 14.95 & 42.50 & -- & \textbf{1.24} \\
H3 & Filter smoothing & \textbf{0.568} & \textbf{1.751} & 4.164 & \textbf{15.69} & 1.157 & \textbf{4.35} & \textbf{2.97} & \textbf{1.85} & -- & 2.78 \\
\midrule
\multicolumn{12}{@{}l}{\emph{\code{periodic}, no loc.}} \\
E10 & ESMDA & 1.575 & \textbf{1.960} & \textbf{4.487} & \textbf{16.19} & 1.602 & -- & 263.32 & \textbf{2.67} & -- & 2.17 \\
F6 & EnKF joint$+$RTPS & 1.194 & 4.084 & 5.071 & 43.68 & 6.091 & \textbf{39.99} & \textbf{23.03} & 43.52 & -- & \textbf{1.26} \\
H10 & Filter smoothing & \textbf{1.127} & 2.632 & 4.621 & 18.45 & \textbf{1.551} & 49.01 & 101.80 & 10.48 & -- & 2.80 \\
\bottomrule
\end{tabular}
\end{adjustbox}
}

\newcommand{\CmpGapTable}{%
\footnotesize
\begin{adjustbox}{max width=\textwidth}
\begin{tabular}{@{}ll lrrr lrrr@{}}
\toprule
case & method & \multicolumn{4}{c}{localization on} & \multicolumn{4}{c}{localization off} \\
\cmidrule(lr){3-6}\cmidrule(lr){7-10}
 & & ID & assim & valid & valid/assim & ID & assim & valid & valid/assim \\
\midrule
\multicolumn{10}{@{}l}{\emph{uDALES truth (matched)}} \\
\code{inflow} (laminar) & ESMDA & E11 & 0.665 & 0.748 & 1.12 & E18 & 0.694 & 0.754 & 1.09 \\
 & EnKF state, no infl. & -- & -- & -- & -- & F14 & 0.236 & 0.570 & 2.41 \\
 & EnKF state$+$RTPS & F8 & 0.167 & 0.626 & 3.75 & F15 & 0.157 & 0.539 & 3.44 \\
 & EnKF joint$+$RTPS & F7 & 0.155 & 0.394 & 2.54 & F13 & 0.144 & 0.319 & 2.21 \\
 & Filter smoothing & H11 & 0.506 & 0.593 & 1.17 & H18 & 0.524 & 0.620 & 1.18 \\
\code{inflow\_turb} & ESMDA & E12 & 0.401 & 0.365 & 0.91 & E19 & 0.456 & 0.407 & 0.89 \\
 & EnKF state, no infl. & -- & -- & -- & -- & F17 & 0.390 & 0.991 & 2.54 \\
 & EnKF state$+$RTPS & F10 & 0.240 & 0.753 & 3.13 & F18 & 0.261 & 0.772 & 2.96 \\
 & EnKF joint$+$RTPS & F9 & 0.206 & 0.559 & 2.71 & F16 & 0.260 & 0.624 & 2.40 \\
 & Filter smoothing & H12 & 0.272 & 0.497 & 1.83 & H19 & 0.256 & 0.441 & 1.72 \\
\code{periodic} & ESMDA & E13 & 1.091 & 1.408 & 1.29 & E20 & 1.163 & 1.671 & 1.44 \\
 & EnKF state, no infl. & -- & -- & -- & -- & F20 & 0.735 & 2.391 & 3.25 \\
 & EnKF state$+$RTPS & F12 & 0.315 & 1.261 & 4.00 & F21 & 0.707 & 2.362 & 3.34 \\
 & EnKF joint$+$RTPS & F11 & 0.317 & 1.551 & 4.89 & F19 & 0.826 & 3.653 & 4.42 \\
 & Filter smoothing & H13 & 0.314 & 1.246 & 3.96 & H20 & 0.706 & 1.766 & 2.50 \\
\midrule
\multicolumn{10}{@{}l}{\emph{PALM truth (model error)}} \\
\code{inflow\_turb} & ESMDA & E2 & 1.422 & 1.377 & 0.97 & E9 & 1.481 & 1.389 & 0.94 \\
 & EnKF joint$+$RTPS & F2 & 0.627 & 1.297 & 2.07 & F5 & 0.880 & 1.891 & 2.15 \\
 & Filter smoothing & H2 & 0.812 & 1.429 & 1.76 & H9 & 0.942 & 1.416 & 1.50 \\
\code{periodic} & ESMDA & E3 & 1.437 & 2.133 & 1.48 & E10 & 1.575 & 1.960 & 1.24 \\
 & EnKF joint$+$RTPS & F3 & 0.704 & 4.072 & 5.78 & F6 & 1.194 & 4.084 & 3.42 \\
 & Filter smoothing & H3 & 0.568 & 1.751 & 3.08 & H10 & 1.127 & 2.632 & 2.34 \\
\bottomrule
\end{tabular}
\end{adjustbox}
}

\newcommand{\CmpObsIntervalTable}{%
\footnotesize
\begin{adjustbox}{max width=\textwidth}
\begin{tabular}{@{}lll rrr rrr rrr@{}}
\toprule
case & method & IDs & \multicolumn{3}{c}{$\alpha$ RMSE [deg]} & \multicolumn{3}{c}{valid.\ RMSE} & \multicolumn{3}{c}{field $\Umag$ RMSE} \\
\cmidrule(lr){4-6}\cmidrule(lr){7-9}\cmidrule(lr){10-12}
 & & (7.5/15/30) & 7.5 & 15 & 30 & 7.5 & 15 & 30 & 7.5 & 15 & 30 \\
\midrule
uDALES, \code{inflow\_turb} & ESMDA & E16/E12/E14 & 2.42 & 2.31 & 2.35 & 0.422 & 0.365 & 0.366 & 0.408 & 0.372 & 0.388 \\
 & Filter smoothing & H16/H12/H14 & 1.94 & 2.61 & 2.86 & 0.471 & 0.497 & 0.487 & 0.539 & 0.533 & 0.548 \\
\addlinespace[2pt]
uDALES, \code{periodic} & ESMDA & E17/E13/E15 & 12.94 & 13.62 & 13.63 & 1.394 & 1.408 & 1.417 & 1.056 & 1.040 & 1.033 \\
 & Filter smoothing & H17/H13/H15 & 13.43 & 14.58 & 14.69 & 1.229 & 1.246 & 1.271 & 1.018 & 1.019 & 1.023 \\
\addlinespace[2pt]
PALM, \code{inflow\_turb} & ESMDA & E6/E2/E4 & 7.83 & 7.12 & 7.44 & 1.327 & 1.377 & 1.323 & 1.024 & 1.082 & 1.068 \\
 & Filter smoothing & H6/H2/H4 & 8.63 & 6.96 & 8.89 & 1.473 & 1.429 & 1.445 & 0.987 & 0.991 & 1.002 \\
\addlinespace[2pt]
PALM, \code{periodic} & ESMDA & E7/E3/E5 & 14.86 & 22.42 & 15.86 & 1.806 & 2.133 & 1.941 & 4.033 & 4.014 & 4.190 \\
 & Filter smoothing & H7/H3/H5 & 14.45 & 15.69 & 18.23 & 1.696 & 1.751 & 1.882 & 4.262 & 4.164 & 4.273 \\
\bottomrule
\end{tabular}
\end{adjustbox}
}

\newcommand{\CmpCostTable}{%
\footnotesize
\begin{adjustbox}{max width=\textwidth}
\begin{tabular}{@{}l lll rrr r@{}}
\toprule
method & analyses & obs per & window & \multicolumn{3}{c}{wall time [h]} & rel.\ to \\
\cmidrule(lr){5-7}
 & (total) & analysis & passes & min & mean & max & EnKF \\
\midrule
ESMDA & 3 (1/window) & 96 & 4 & 1.77 & 2.15 & 2.36 & 1.77$\times$ \\
EnKF state, no infl. & 180 (60/window) & 12 & 1 & 1.11 & 1.17 & 1.20 & 0.96$\times$ \\
EnKF state$+$RTPS & 180 (60/window) & 12 & 1 & 1.01 & 1.17 & 1.26 & 0.96$\times$ \\
EnKF joint$+$RTPS & 180 (60/window) & 12 & 1 & 1.06 & 1.21 & 1.28 & 1.00$\times$ \\
Filter smoothing & 183 (61/window) & 12 / 96 & $\approx 4$ & 2.35 & 2.77 & 2.96 & 2.28$\times$ \\
\bottomrule
\end{tabular}
\end{adjustbox}
}

% --- numbers quoted in the prose --------------------------------
\newcommand{\CmpFlagAngleE}{2.31}
\newcommand{\CmpFlagAngleJ}{3.36}
\newcommand{\CmpFlagAngleH}{2.61}
\newcommand{\CmpFlagValidE}{0.365}
\newcommand{\CmpFlagValidJ}{0.559}
\newcommand{\CmpFlagValidH}{0.497}
\newcommand{\CmpPerAssimE}{1.091}
\newcommand{\CmpPerAssimH}{0.314}
\newcommand{\CmpPerVmagJ}{4.221}
\newcommand{\CmpPerVmagH}{0.756}



\subsection{What is compared --- and what cannot be}
\label{sec:cmp:scope}

Sections~\ref{sec:esmda-takeaways}, \ref{sec:filt-takeaways} and
\ref{sec:fs:takeaways} each rank variants \emph{within} one campaign. This
section puts the three campaigns on one axis. The comparison is organised
around \emph{matched cells}: a cell is one
(truth model, boundary-condition case, localization) triple, and every method
that was run at those settings supplies one row of it. Five method variants
appear:

\begin{description}[nosep,leftmargin=2.2em,style=nextline]
  \item[ESMDA (IDs \code{E*})] the parameter smoother of \S\ref{sec:esmda}; one
        analysis per window, three MDA steps, state never updated
        (\code{obs15} runs only).
  \item[EnKF state, no infl.\ (IDs \code{F*})] a stochastic EnKF updating the
        state every \SI{2}{\second} with no inflation.
  \item[EnKF state$+$RTPS (IDs \code{F*})] the same with RTPS inflation
        ($\alpha = 0.6$).
  \item[EnKF joint$+$RTPS (IDs \code{F*})] the same filter with the parameters
        appended to the state vector.
  \item[Filter smoothing (IDs \code{H*})] the hybrid of \S\ref{sec:fs}: ESMDA
        parameters $+$ EnKF state inside every window (\code{obs15} runs only).
\end{description}

Every run keeps the ID it was given in Sections~\ref{sec:esmda}--\ref{sec:fs};
the ID columns of Tables~\ref{tab:cmp:master-ud}--\ref{tab:cmp:obsinterval} are
the map back to those sections, and the folder names are listed only in the
experiment-matrix table of each method section.

That gives $45$ matched cells; $39$ produced a \code{run\_summary.yaml} and $6$
did not (see \S\ref{sec:cmp:caveats}). The observation interval is held at
\SI{15}{\second} wherever a campaign sweeps it, so the master tables are a
clean method contrast; the sweep itself is compared separately in
Table~\ref{tab:cmp:obsinterval}.

\paragraph{Four things that are \emph{not} comparable across campaigns.}
These caveats are load-bearing: several apparent method differences below are
partly definitional, and the text flags them where they bite.

\begin{enumerate}[nosep,leftmargin=1.6em]
  \item \textbf{The observation stream is not the same.} ESMDA assimilates
        \emph{aggregated} observations --- $96$ values per window at
        $\Delta t_{\mathrm{obs}} = \SI{15}{\second}$ (6 sensors $\times$ 2
        components $\times$ 8 bins). Both filters assimilate the \emph{raw}
        stream, $12$ values every \SI{2}{\second}, i.e.\ $720$ values per
        window. The hybrid does both. ESMDA is therefore working from less
        data by construction, and any ESMDA win below is a win in spite of
        that.
  \item \textbf{``Reduction vs.\ prior'' has three different meanings}, so
        only \emph{absolute} parameter RMSE is quoted here. For ESMDA and the
        hybrid the baseline is that run's own re-drawn climatological prior
        ($14.5$--$19.0^\circ$ on the angle). For joint filtering it is the
        random-walk-widened \emph{forecast} prior ($\approx 29$--$30^\circ$).
        For state-only filtering the parameter ensemble is never touched, so
        the reduction is exactly zero and the reported RMSE is the static
        prior ($19.29^\circ$, $0.919$~m/s) --- those entries are printed in
        italics with a dagger and are excluded from the ``best'' bolding.
  \item \textbf{The observation-interval knob means different things.} In the
        ESMDA campaign \code{obs7p5/15/30} changes the number of assimilated
        observations ($192/96/48$ per window). In the filter-smoothing
        campaign it changes only the \emph{smoother's} aggregation bin; the
        filter still assimilates every \SI{2}{\second} regardless. The two
        sweeps are not the same experiment.
  \item \textbf{Calibration is measured with different instruments.} The
        filters and the hybrid report an innovation $\chi^2$; ESMDA has no
        sequential innovation and reports a normalised data mismatch $O_N$
        instead, so its $\chi^2$ column is empty everywhere. The one
        diagnostic all five share is the $z$-score standard deviation
        (parameters and validation sensors), and that is what the calibration
        discussion leans on.
\end{enumerate}

Everything else follows the conventions of \S\ref{sec:setup}: sensor RMSEs are
the 3-component \emph{velocity-vector} RMSE
(\code{sensor\_metrics.*.velocity\_vector\_rmse.mean}), \emph{field $\Umag$} is
\code{state\_metrics.vel\_magnitude\_rmse.mean}, $q$ is the field hit rate, and
figure legends print a magnitude-only $\Umag$ error that is roughly half the
tabulated vector RMSE.

\subsection{Master comparison tables}
\label{sec:cmp:master}

\begin{table}[H]
  \centering
  \caption{Matched-model comparison (uDALES truth). Rows are method variants
  within one (case, localization) cell; bold marks the best value in each
  column \emph{within that cell}. $z_{\mathrm{par}}$ is the pooled parameter
  $z$-score std and $z_{\mathrm{val}}$ the validation-sensor
  $\Umag$ $z$-score std (both $1.03$ when calibrated); $q$ is the field hit
  rate. $\dagger$ marks carried-prior entries: state-only filters never update
  the parameter ensemble, so those numbers are the static prior, not a
  posterior. Two readings deserve care --- the \emph{assim} column rewards
  methods that update state directly at the sensors and is a fit diagnostic
  rather than a skill measure (\S\ref{sec:cmp:generalisation}), and a high $q$
  can be bought with a diverging ensemble (\S\ref{sec:cmp:calibration}).}
  \label{tab:cmp:master-ud}
  \CmpMasterUDALES
\end{table}

\begin{table}[H]
  \centering
  \caption{Model-error comparison (PALM truth, uDALES assimilation model).
  Same layout as Table~\ref{tab:cmp:master-ud}; \code{field\_metrics} is not
  written for PALM-truth runs, hence the empty $q$ column. Both laminar
  \code{inflow} cells are empty because the PALM truth generation itself
  diverged in all three campaigns.}
  \label{tab:cmp:master-palm}
  \CmpMasterPALM
\end{table}

\begin{table}[H]
  \centering
  \caption{Fit versus generalisation. \emph{valid/assim} is the ratio of
  held-out to assimilated velocity-vector RMSE: a value near $1$ means the
  method is as good away from the sensors as at them, and a large value means
  the fit is local. This is the sharpest structural difference between the
  parameter smoother and the state filters.}
  \label{tab:cmp:gap}
  \CmpGapTable
\end{table}

\begin{table}[H]
  \centering
  \caption{Observation-interval sweeps, ESMDA versus filter smoothing
  (correlation localization throughout). Recall caveat~3 of
  \S\ref{sec:cmp:scope}: the same column heading denotes a change in the
  assimilated observation count for ESMDA but only a change in the smoother's
  aggregation bin for the hybrid. The sequential-filtering campaign has no
  such sweep.}
  \label{tab:cmp:obsinterval}
  \CmpObsIntervalTable
\end{table}

\begin{table}[H]
  \centering
  \caption{Cost. \emph{analyses} counts Kalman/MDA updates over the whole
  \SI{360}{\second} run; \emph{obs per analysis} is the number of scalar
  observations each update sees; \emph{window passes} is how many times the
  \SI{120}{\second} window is integrated by the full ensemble per assimilation
  window. Wall-time statistics are over all completed runs of that method
  (both truth models and all cases), so they mix solver speeds --- the
  periodic uDALES runs are intrinsically the fastest. The last column
  normalises by the joint-EnKF mean.}
  \label{tab:cmp:cost}
  \CmpCostTable
\end{table}

\subsection{Summary figures}
\label{sec:cmp:figures}

\subsubsection{State at the validation sensors}
\label{sec:cmp:fig-state}

Figures~\ref{fig:cmp:valid-ud-inflow}--\ref{fig:cmp:valid-noloc} put the
per-run sensor figures of Sections~\ref{sec:esmda}--\ref{sec:fs} side by side.
Every panel triple is the same crop: ground sensor $0$ at $(2,30,2)$, the
elevated sensor $3$ at $(10,30,20)$ --- the held-out sensor furthest from any
observation --- and the per-cycle $\Umag$ error. Two reading rules apply
throughout. \emph{(i)} The filter-smoothing column is plotted against
\emph{cycle index} ($0$--$180$, \SI{2}{\second} per cycle) while the ESMDA and
EnKF columns are plotted against seconds ($0$--$360$); the physical duration is
identical and the curves are not redrawn. \emph{(ii)} The panel legends print a
magnitude-only $\Umag$ RMSE, roughly half the velocity-\emph{vector} RMSE
quoted in the label strips and in the tables; the two must never be compared
directly.

\begin{figure}[H]
  \centering
  \includegraphics[width=\textwidth]{comparison/cmp_valid_ud_inflow.png}
  \caption{Held-out sensors, uDALES truth, laminar \code{inflow}, localization
  on: E11 (ESMDA), F7 (EnKF joint$+$RTPS), H11 (filter smoothing). The joint
  filter (middle) is the only column whose ensemble stays wide enough to
  bracket the truth; both smoother-based columns collapse onto a narrow,
  phase-shifted mean --- the laminar case where the boundary parameters cannot
  express the interior error. \emph{x-axis:} seconds for E11/F7, cycle index
  for H11.}
  \label{fig:cmp:valid-ud-inflow}
\end{figure}

\begin{figure}[H]
  \centering
  \includegraphics[width=\textwidth]{comparison/cmp_valid_ud_inflow_turb.png}
  \caption{Held-out sensors, uDALES truth, \code{inflow\_turb}, localization
  on: E12, F9, H12. This is the case where the ordering reverses. E12 tracks
  the elevated sensor with a narrow, slightly low-biased ensemble; F9 has a
  visibly wider ensemble that brackets the truth but whose mean is flatter;
  H12 sits between them. \emph{x-axis:} seconds for E12/F9, cycle index for
  H12.}
  \label{fig:cmp:valid-ud-turb}
\end{figure}

\begin{figure}[H]
  \centering
  \includegraphics[width=\textwidth]{comparison/cmp_valid_ud_periodic.png}
  \caption{Held-out sensors, uDALES truth, \code{periodic}, localization on:
  E13, F11, H13. All three columns show the spin-up transient of the periodic
  configuration and a mean that never recovers the truth amplitude; F11's
  ensemble is the widest, which is what buys it the highest hit rate $q$
  despite the worst field RMSE (\S\ref{sec:cmp:calibration}).
  \emph{x-axis:} seconds for E13/F11, cycle index for H13.}
  \label{fig:cmp:valid-ud-periodic}
\end{figure}

\begin{figure}[H]
  \centering
  \includegraphics[width=\textwidth]{comparison/cmp_valid_ud_stateonly.png}
  \caption{The state-only companions of
  Figures~\ref{fig:cmp:valid-ud-inflow}--\ref{fig:cmp:valid-ud-periodic}:
  F8 (laminar \code{inflow}), F10 (\code{inflow\_turb}) and F12
  (\code{periodic}), all EnKF state$+$RTPS with localization on. They are shown
  separately because a fourth column would push the axis labels below legibility
  at \code{\textbackslash textwidth}. The ensembles are the widest of any
  variant --- RTPS is doing its job --- but the mean carries no
  boundary-condition information, so the elevated sensor is tracked worst of
  all here. All three x-axes are in seconds.}
  \label{fig:cmp:valid-stateonly}
\end{figure}

\begin{figure}[H]
  \centering
  \includegraphics[width=\textwidth]{comparison/cmp_valid_palm_inflow_turb.png}
  \caption{Model error: held-out sensors, PALM truth, \code{inflow\_turb},
  localization on: E2, F2, H2. Compare with
  Figure~\ref{fig:cmp:valid-ud-turb}, the matched-model version of the same
  cell. The truth trace is visibly outside the ensemble for long stretches in
  all three columns --- the state floor of \S\ref{sec:cmp:modelerror} at a
  single sensor. \emph{x-axis:} seconds for E2/F2, cycle index for H2.}
  \label{fig:cmp:valid-palm-turb}
\end{figure}

\begin{figure}[H]
  \centering
  \includegraphics[width=\textwidth]{comparison/cmp_valid_palm_periodic.png}
  \caption{Model error: held-out sensors, PALM truth, \code{periodic},
  localization on: E3, F3, H3. F3 is the divergent run of
  Table~\ref{tab:cmp:master-palm} ($\alpha$ RMSE $138.50^\circ$): its ensemble
  overshoots the truth by a factor $\approx 3$ from $t \approx \SI{40}{\second}$
  onwards and never returns. H3 stays bounded.
  \emph{x-axis:} seconds for E3/F3, cycle index for H3.}
  \label{fig:cmp:valid-palm-periodic}
\end{figure}

\begin{figure}[H]
  \centering
  \includegraphics[width=\textwidth]{comparison/cmp_valid_ud_noloc.png}
  \caption{Localization off, uDALES truth: \code{inflow\_turb} (top row: E19,
  F16, H19) and \code{periodic} (bottom row: E20, F19, H20). The counterparts
  with localization on are Figures~\ref{fig:cmp:valid-ud-turb}
  and~\ref{fig:cmp:valid-ud-periodic}. On \code{inflow\_turb} switching
  localization off costs little; on \code{periodic} it is the difference
  between a tracked and an untracked sensor for every method.
  \emph{x-axis:} seconds for the E and F columns, cycle index for the H
  column.}
  \label{fig:cmp:valid-noloc}
\end{figure}

\subsubsection{Parameter estimation}
\label{sec:cmp:fig-params}

Figures~\ref{fig:cmp:paramerr-ud}--\ref{fig:cmp:marg-palm} show the same three
parameter-estimating methods on the same axes. The state-only EnKF runs are
omitted: their parameter ensemble is never updated, so every panel would be the
flat static prior ($19.29^\circ$, $0.919$~m/s). Unlike the sensor figures, the
parameter-error panels are in seconds for \emph{all} three methods. The
marginals carry one asymmetry that must be read carefully: ESMDA and the hybrid
resolve the boundary condition on a $15$-knot spline and their ``final knot''
is knot $14$ of $15$, whereas the joint EnKF carries only $3$ knots in its
state vector and its final knot is knot $2$ of $3$ --- the panels are the same
quantity at a coarser temporal resolution.

\begin{figure}[H]
  \centering
  \includegraphics[width=\textwidth]{comparison/cmp_param_error_ud.png}
  \caption{Parameter-error time series, uDALES truth, localization on: inflow
  angle $\alpha$ (upper panel of each pair) and velocity magnitude $\Umag$
  (lower). Rows are the three cases; columns are ESMDA (E11/E12/E13), EnKF
  joint$+$RTPS (F7/F9/F11) and filter smoothing (H11/H12/H13). The
  smoother-based columns step once per assimilation window and the filter
  column updates continuously, which is why only the middle column is a rough
  trace. The $\Umag$ row of F7/F9 is the clearest single view of the result of
  \S\ref{sec:cmp:params}: the joint filter's magnitude error never leaves the
  prior band.}
  \label{fig:cmp:paramerr-ud}
\end{figure}

\begin{figure}[H]
  \centering
  \includegraphics[width=\textwidth]{comparison/cmp_param_error_palm.png}
  \caption{Parameter-error time series under model error (PALM truth,
  localization on): \code{inflow\_turb} (E2, F2, H2) and \code{periodic}
  (E3, F3, H3). Same layout as Figure~\ref{fig:cmp:paramerr-ud}. The
  \code{periodic} row of F3 shows the divergence directly: the angle error
  saturates near $150^\circ$ within \SI{40}{\second} and the magnitude error
  near $20$~m/s, while E3 and H3 stay in the $10$--$30^\circ$ band.}
  \label{fig:cmp:paramerr-palm}
\end{figure}

\begin{figure}[H]
  \centering
  \includegraphics[width=\textwidth]{comparison/cmp_marginals_ud.png}
  \caption{Final-knot prior (grey) and posterior (blue) marginals, uDALES
  truth, localization on; dashed line is the truth. Same rows and columns as
  Figure~\ref{fig:cmp:paramerr-ud}. On both inflow cases every posterior has
  contracted sharply, and on \code{inflow\_turb} the ESMDA and hybrid angle
  posteriors sit on the truth line while their $\Umag$ posteriors sit visibly
  above it ($z = -4.60$ and $-6.28$) --- contracted \emph{and} displaced, the
  over-confidence of \S\ref{sec:cmp:calibration}. On \code{periodic} the
  posteriors have barely contracted and have moved away from the truth.}
  \label{fig:cmp:marg-ud}
\end{figure}

\begin{figure}[H]
  \centering
  \includegraphics[width=\textwidth]{comparison/cmp_marginals_palm.png}
  \caption{Final-knot marginals under model error (PALM truth, localization
  on), same layout as Figure~\ref{fig:cmp:paramerr-palm}. The posteriors are as
  narrow as in the matched case but are displaced from the truth line: the
  ``confidently wrong'' behaviour of \S\ref{sec:cmp:modelerror}. F3's angle
  posterior has left the plotted prior range entirely.}
  \label{fig:cmp:marg-palm}
\end{figure}

\subsubsection{Aggregate metrics}
\label{sec:cmp:fig-aggregate}

\begin{figure}[H]
  \centering
  \includegraphics[width=\textwidth]{comparison/cmp_accuracy.png}
  \caption{Held-out-sensor velocity-vector RMSE (left) and domain field
  $\Umag$ RMSE (right) for every matched cell. Missing bars are
  configurations that were not run (the uninflated state-only filter exists
  only at localization off, and no state-only variant was run for PALM truth).
  Note that the ordering flips between cases: EnKF joint$+$RTPS wins on
  laminar \code{inflow}, ESMDA on \code{inflow\_turb}, and filter smoothing on
  \code{periodic}.}
  \label{fig:cmp:accuracy}
\end{figure}

\begin{figure}[H]
  \centering
  \includegraphics[width=\textwidth]{comparison/cmp_parameters.png}
  \caption{Parameter recovery, plotted as lollipops \emph{from the static
  climatological prior} (dashed line, $19.29^\circ$ and $0.919$~m/s) on a
  logarithmic axis --- a stem pointing down is an improvement on climatology,
  one pointing up is a degradation. Only the three parameter-estimating
  methods appear; state-only filters sit exactly on the dashed line by
  construction. The inflow angle is recoverable on \code{inflow} and
  \code{inflow\_turb} by all three; the velocity magnitude is only improved by
  the two \emph{smoother}-based methods, while the joint filter leaves it at
  the prior in the matched cases and drives it far past the prior on
  \code{periodic} and under model error.}
  \label{fig:cmp:parameters}
\end{figure}

\begin{figure}[H]
  \centering
  \includegraphics[width=\textwidth]{comparison/cmp_calibration.png}
  \caption{Calibration on a logarithmic axis, all matched cells (UD $=$ uDALES
  truth, PA $=$ PALM truth). Dashed lines mark the calibrated values
  ($\chi^2 = 1$; $\mathrm{std}(z) = 1.03$ for $N = 50$). ESMDA is absent
  from the top panel because it reports no innovation $\chi^2$
  (caveat~4, \S\ref{sec:cmp:scope}); state-only filters are absent from the
  bottom panel because their parameter $z$ is a prior, not a posterior. The
  RTPS-inflated filters are the only variants close to calibrated, and they
  are precisely the variants with the least parameter skill.}
  \label{fig:cmp:calibration}
\end{figure}

\begin{figure}[H]
  \centering
  \includegraphics[width=\textwidth]{comparison/cmp_generalisation.png}
  \caption{Fit versus generalisation: assimilated-sensor RMSE (horizontal)
  against held-out-sensor RMSE (vertical), log--log, one marker per completed
  run. Points on the dashed $1{:}1$ line generalise as well as they fit;
  points far above it fit the six assimilated sensors much better than the
  four held-out ones. ESMDA (blue circles) sits on or below the line on
  \code{inflow\_turb}; every filter variant sits well above it. Black marker
  edges denote correlation localization.}
  \label{fig:cmp:generalisation}
\end{figure}

\subsection{Discussion}
\label{sec:cmp:discussion}

\subsubsection{State accuracy: the winner changes with the case, not with the
method}
\label{sec:cmp:state}

No method is best everywhere, and the ranking is not even stable in sign.
Figure~\ref{fig:cmp:accuracy} is the whole result in one panel pair; reading
held-out-sensor RMSE off Table~\ref{tab:cmp:master-ud} (localization on,
matched model):

\begin{itemize}[nosep]
  \item \code{inflow} (laminar): joint EnKF F7 $0.394$ $<$ hybrid H11 $0.593$
        $<$ EnKF state$+$RTPS F8 $0.626$ $<$ ESMDA E11 $0.748$. The filter wins
        by $47\%$ over the smoother.
  \item \code{inflow\_turb}: ESMDA E12 $0.365$ $<$ hybrid H12 $0.497$ $<$ joint
        EnKF F9 $0.559$ $<$ EnKF state$+$RTPS F10 $0.753$. The smoother wins by
        $35\%$ over the joint filter --- the ordering has \emph{reversed}.
  \item \code{periodic}: hybrid H13 $1.246$ $\approx$ EnKF state$+$RTPS F12
        $1.261$ $<$ ESMDA E13 $1.408$ $<$ joint EnKF F11 $1.551$.
\end{itemize}

The domain field tells the same story with the same reversals: field $\Umag$
RMSE is $0.247$ (F7) on laminar, $0.372$ (E12) on \code{inflow\_turb} and
$1.019$ (H13) on \code{periodic}. Figures~\ref{fig:cmp:valid-ud-inflow},
\ref{fig:cmp:valid-ud-turb} and~\ref{fig:cmp:valid-ud-periodic} show the same
three orderings at the held-out sensors themselves.

The mechanism is the one the individual sections each identified from their
own side. Where the boundary forcing is the dominant control on the interior
flow --- the turbulent inlet --- correcting the forcing propagates everywhere
and beats correcting local state. Where the interior flow is dominated by a
phase error the boundary parameters cannot express (laminar inflow, whose
parameters ESMDA cannot pin down: \S\ref{sec:esmda}) or by internal dynamics
that are not driven by the inlet at all (periodic lateral boundaries), a state
filter that simply nudges the field wins or ties. \good{The practical
consequence is that the choice between a parameter smoother and a state filter
is a statement about which error source dominates, not a statement about which
algorithm is better.}

One number in Table~\ref{tab:cmp:master-ud} should \emph{not} be read as skill:
the assimilated-sensor RMSE. The filters update the state directly at those six
sensors, so they beat ESMDA in that column in all $10$ cells where
they compete --- by a factor $1.9$ on \code{inflow\_turb} (F9 $0.206$ vs.\
E12 $0.401$) and $3.5$ on \code{periodic} (F12 $0.315$ vs.\ E13 $1.091$) ---
while losing
the held-out column in the first of those. It is a fit diagnostic;
\S\ref{sec:cmp:generalisation} treats the gap as the quantity of interest.

\subsubsection{Parameter identifiability is a property of the case, and of
\emph{which} parameter}
\label{sec:cmp:params}

Figure~\ref{fig:cmp:parameters} separates two claims that are easy to conflate.

\paragraph{The inflow angle is recoverable on both inflow cases, by all three
parameter-estimating methods.} On \code{inflow\_turb} with localization the
angle RMSE is $2.31^\circ$ (E12), $2.61^\circ$ (H12) and $3.36^\circ$ (F9)
against a $\approx 17$--$19^\circ$ climatological prior; without localization,
$2.45$ / $3.45$ / $4.88^\circ$ (E19 / H19 / F16). On laminar \code{inflow} the
ordering inverts --- joint EnKF $3.90^\circ$ (F13, no loc) and $5.26^\circ$
(F7, loc) beat the hybrid ($7.87$ / $5.81^\circ$) and ESMDA ($7.59$ /
$7.74^\circ$)
--- consistent with the laminar case being where the smoother stalls on phase
drift while the filter still extracts directional information from the raw
\SI{2}{\second} stream.

\paragraph{The velocity magnitude is only recovered by the smoother-based
methods.} Against the $0.919$~m/s static prior, ESMDA E12 reaches $0.369$~m/s
and the hybrid H12 $0.433$~m/s on \code{inflow\_turb} with localization
($0.439$ and $0.365$ without, E19 and H19), while the joint EnKF reports
$0.912$ (F9) and $0.955$~m/s (F16) --- i.e.\ no improvement at all, and a
slight degradation without localization. The lower row of each column pair in
Figure~\ref{fig:cmp:paramerr-ud} shows this directly: F7 and F9 never leave the
prior band. The tempered, three-step MDA update extracts the magnitude; the
single linear Kalman step appended to a state vector does not. \bad{Worse, the
joint filter's magnitude update is actively unsafe where the parameter is
unidentifiable}: on \code{periodic} it reaches $4.221$~m/s (F11, matched) and
$21.276$~m/s (F3, PALM truth), against $0.756$ (H13) and $1.157$~m/s (H3) for
the hybrid at the same settings. The unbounded Gaussian update drives members
to physically impossible inflow speeds, and the state follows (field $\Umag$
RMSE $1.733$ vs.\ $1.019$).

\paragraph{On \code{periodic} nothing works, and that is a property of the
observation network.} All three parameter-estimating methods (E13, F11, H13)
land within $13.6$--$19.0^\circ$ of a $\approx 16.6^\circ$ prior; the bottom
row of Figure~\ref{fig:cmp:marg-ud} shows why --- the posteriors have barely
contracted. Six near-ground sensors
under periodic lateral boundaries do not constrain an inflow direction that no
longer drives the flow. The correct conclusion is not ``use a different
estimator'' but ``do not estimate this parameter in this configuration''.

\subsubsection{Generalisation to held-out and elevated sensors}
\label{sec:cmp:generalisation}

Table~\ref{tab:cmp:gap} and Figure~\ref{fig:cmp:generalisation} isolate the one
difference that is structural rather than case-dependent.

\begin{itemize}[nosep]
  \item \textbf{ESMDA generalises essentially perfectly}: valid/assim
        $= 0.89$--$1.48$ over all ten matched cells, and \emph{below one} on
        both \code{inflow\_turb} cells (E12 $0.91$, E19 $0.89$) and both PALM
        \code{inflow\_turb} cells (E2 $0.97$, E9 $0.94$) --- it is more
        accurate at
        the four sensors it never saw than at the six it fitted. A global
        boundary-condition correction has no local overfitting mode.
  \item \textbf{Every filter variant fits locally}: valid/assim
        $= 2.2$--$4.9$ on matched-model cells, i.e.\ held-out error is two to
        five times the fitted error. RTPS makes this \emph{worse}, not better
        (laminar, no loc: F14 $2.41 \rightarrow$ F15 $3.44$), because inflation
        buys
        calibration at the sensors without adding any information away from
        them.
  \item \textbf{The hybrid inherits the smoother's behaviour on the inflow
        cases and the filter's on periodic}: $1.17$--$1.83$ on \code{inflow}
        and \code{inflow\_turb}, but $3.96$ (H13) on \code{periodic} with
        localization --- exactly the case where its parameter half has nothing
        to contribute and only the filter half is doing work.
\end{itemize}

Figure~\ref{fig:cmp:valid-ud-turb} shows what this looks like at the elevated
$z = \SI{20}{\meter}$ validation sensor, the one furthest from any
observation. All three methods track it; the difference is in the ensemble.
E12 produces a narrow ensemble whose mean follows the truth but sits
systematically below it in the second half; F9 produces a much wider ensemble
that brackets truth but whose mean is flatter; H12 is intermediate. The
state-only companion F10 (Figure~\ref{fig:cmp:valid-stateonly}) is the extreme
of the same trade: the widest ensemble of all and the worst held-out mean. That is the accuracy/calibration trade-off of
\S\ref{sec:cmp:calibration} made visible at a single sensor.

\subsubsection{Calibration: the well-calibrated methods are the ones with no
parameter skill}
\label{sec:cmp:calibration}

This is the least comfortable result of the comparison
(Figure~\ref{fig:cmp:calibration}).

\begin{itemize}[nosep]
  \item \textbf{Innovation $\chi^2$}: the RTPS-inflated filters are close to
        calibrated on five of the six matched cells ($0.92$--$1.82$), the one
        exception being \code{periodic} without localization (F21, $14.05$).
        The hybrid, whose filter half is the \emph{same} EnKF with the
        \emph{same} RTPS, reports $13.26$ / $14.96$ on laminar (H11 / H18),
        $3.68$ / $3.43$ on \code{inflow\_turb} (H12 / H19) and $1.81$ / $23.06$
        on \code{periodic} (H13 / H20): the
        smoother's parameter update injects an error that the filter then has
        to absorb, and it shows up as an inflated innovation. Removing
        inflation is worse still ($3.98$, $8.79$, $27.72$).
  \item \textbf{Validation-sensor $z$}: same ordering. The RTPS-inflated filters are
        honest ($0.83$--$1.48$ on matched cells except periodic without
        localization; the uninflated filter F20 reaches $7.00$), ESMDA
        is over-confident ($1.32$--$7.82$), the hybrid most of all
        ($1.40$--$11.25$). Figure~\ref{fig:cmp:marg-ud} shows the same
        over-confidence on the parameters: the E and H posteriors are the
        narrowest and the furthest from the truth line.
  \item \textbf{Parameter $z$}: no method is calibrated. The joint EnKF is
        closest ($0.74$--$3.29$ on five of six matched cells, $11.54$ for F19
        on the sixth) but, as
        \S\ref{sec:cmp:params} showed, its parameter ensemble is wide precisely
        because it is not learning the magnitude. ESMDA spans
        $2.58$--$89.85$ and the hybrid $2.67$--$41.51$; without localization
        both can reach three digits ($263.32$ for E10 on PALM
        \code{periodic}, $101.80$ for H10).
  \item \textbf{Hit rate $q$ is not a skill measure.} The highest matched-model
        $q$ in the whole set, $0.887$, belongs to F11 (joint EnKF,
        \code{periodic}) --- the same run with field $\Umag$ RMSE $1.733$ and a
        $4.221$~m/s parameter error. A diverging ensemble covers the truth.
        Read $q$ next to a RMSE column, never alone.
\end{itemize}

The common cause is the same one every section flagged individually: every run
uses $\sigma_{\mathrm{obs}} = 0.1$~m/s with
\code{caveat: no\_representativeness\_error}. An observation error that
excludes model representativeness makes the analysis too confident in all
three algorithms; it is a configuration property, not a property of any one
method.

\subsubsection{Robustness to model error}
\label{sec:cmp:modelerror}

Switching the truth model to PALM costs a factor $3$--$4$ in held-out RMSE for
every method, and --- more importantly --- \emph{compresses the methods
together}. On \code{inflow\_turb} with localization the matched-model spread
is $0.365$--$0.753$ (a factor $2.1$ between best and worst); under model error
it is $1.297$--$1.429$ (a factor $1.10$). The ordering also changes: F2
($1.297$) now edges out E2 ($1.377$) and H2 ($1.429$), a reversal of the
matched result. Figure~\ref{fig:cmp:valid-palm-turb} is the sensor-level view
of that compression: the three columns are barely distinguishable. \bad{Once representativeness error dominates,
the choice of assimilation algorithm stops mattering for state accuracy.}

Parameter estimates do not degrade gracefully --- they become
\emph{confidently wrong}. On PALM \code{inflow\_turb} all three methods report
$6.6$--$11.1^\circ$ angle RMSE, which looks like partial success, but with
pooled parameter $z$ of $12.49$ (E2) and $12.44$ (H2); E2's data mismatch
stalls at $O_N = 30.1$ instead of descending to E12's matched-model $2.37$ (and
at $57.4$ on PALM \code{periodic}). The posterior is narrow and displaced ---
Figure~\ref{fig:cmp:marg-palm}.

Where the methods do separate under model error is stability. On PALM
\code{periodic} with localization the joint EnKF F3 diverges outright --- angle
RMSE $138.50^\circ$, magnitude $21.276$~m/s, held-out RMSE $4.072$, validation
$z$ of $42.50$ --- while the hybrid H3 stays bounded ($15.69^\circ$,
$1.157$~m/s, held-out $1.751$, $z = 1.85$) and ESMDA E3 lands in between
($22.42^\circ$, $1.090$~m/s, $2.133$). Figures~\ref{fig:cmp:valid-palm-periodic}
and~\ref{fig:cmp:paramerr-palm} show the divergence itself. \good{Splitting the parameter update off
into a tempered smoother is what prevents the blow-up}: three damped MDA steps
against $96$ aggregated observations cannot move the parameters as far in one
window as $60$ untempered Kalman updates can. That robustness, not accuracy, is
the strongest argument for the hybrid in this set.

The field is beyond repair in the periodic model-error corner regardless of
method: field $\Umag$ RMSE is $4.01$--$5.91$ against $\approx 1.02$ matched.

\subsubsection{Cost}
\label{sec:cmp:cost}

Table~\ref{tab:cmp:cost} gives the accounting. Per \SI{360}{\second} run on 10
parallel processes, the sequential filters cost $1.01$--$1.28$~h, ESMDA
$1.77$--$2.36$~h ($1.77\times$ the joint-filter mean) and filter smoothing
$2.35$--$2.96$~h ($2.28\times$). The driver is the number of times the ensemble
must re-integrate each window: one pass for a filter, four for ESMDA (prior
$+$ three MDA steps), and effectively four for the hybrid, which runs both.

Cost and accuracy are not aligned, and they are not anti-aligned either:

\begin{itemize}[nosep]
  \item On laminar \code{inflow} the \emph{cheapest} method is also the best
        on every accuracy column: joint EnKF F13, held-out $0.319$ at
        $1.19$~h. There is no case for paying $2.4\times$ for the hybrid here.
  \item On \code{inflow\_turb} the best held-out RMSE ($0.365$) costs
        $2.30$~h with ESMDA E12, against $0.559$ at $1.27$~h with the joint
        filter F9 --- a $35\%$ accuracy gain for an $81\%$ cost increase, plus
        the parameter estimate.
  \item On \code{periodic} the hybrid H13 ($1.246$ at $2.38$~h) is
        statistically indistinguishable from EnKF state$+$RTPS F12 ($1.261$ at
        $1.01$~h) on the state. The extra $1.4$~h buys only a bounded
        parameter estimate --- worth it if the boundary condition is a
        deliverable, not otherwise.
\end{itemize}

The wall-time ranges in Table~\ref{tab:cmp:cost} mix cases and truth models and
are therefore contaminated by solver speed (uDALES is fastest on the periodic
configuration); the within-cell comparisons above are the reliable ones.

\subsubsection{Which method for which regime}
\label{sec:cmp:recommend}

\begin{description}[leftmargin=2.4em,style=nextline]
  \item[Matched model, forcing-dominated flow (\code{inflow\_turb}) ---
        \good{ESMDA} (E12).] Best held-out RMSE ($0.365$), best field
        ($0.372$), best parameters ($2.31^\circ$, $0.369$~m/s), generalises
        better than it fits, and cheaper than the hybrid. Accept that it is
        over-confident ($z_{\mathrm{val}} = 2.71$) --- if the ensemble spread
        is a deliverable, take H12 instead and pay $27\%$ more wall time for
        $z_{\mathrm{val}} = 2.36$ and $0.497$ held-out.
  \item[Matched model, laminar inflow --- \good{EnKF joint$+$RTPS} (F13).]
        Best on every accuracy column, best calibrated ($\chi^2 =
        0.90$--$0.99$, $z_{\mathrm{val}} = 0.83$--$0.88$), cheapest, and it
        recovers the angle ($3.90^\circ$) which the smoother cannot here. Do not expect the
        magnitude ($0.818$--$0.910$~m/s, essentially the prior). Run it
        \emph{without} localization: localization costs $24\%$ of held-out
        RMSE on this case. Figure~\ref{fig:cmp:valid-noloc} shows the
        localization-off cells of the other two cases, where the sign of the
        effect is the opposite.
  \item[Matched model, periodic --- \good{EnKF state$+$RTPS} (F12), or
        \good{filter smoothing} (H13) if a boundary estimate is required.] They
        tie on the state ($1.261$ vs.\ $1.246$ held-out; $1.029$ vs.\ $1.019$
        field). \bad{Do not use the joint EnKF here} (F11): it degrades both
        the state ($1.733$ field) and the parameters ($4.221$~m/s). Use
        localization: it is worth $47$--$58\%$ of held-out RMSE on this case.
  \item[Model error (PALM truth) --- \good{ESMDA or filter smoothing,
        with localization; parameters not trustworthy}.] All methods collapse
        onto a $1.30$--$1.43$ held-out floor on \code{inflow\_turb}, so pick on
        stability rather than accuracy. \bad{Avoid the joint EnKF on periodic}
        (F3, divergence). Localization is mandatory: without it the hybrid's
        pooled parameter $z$ reaches $101.80$ (H10) and ESMDA's $263.32$
        (E10).
  \item[If you need both a state and a boundary-condition estimate and cannot
        tune per case --- \good{filter smoothing}.] Scoring each method by its
        \emph{regret} --- how much worse its held-out RMSE is than the best
        method in the same cell --- gives, over the ten matched cells: hybrid
        median $+10.2\%$, worst case $+94\%$; ESMDA median $+6.2\%$,
        worst case $+136\%$; joint EnKF median $+53.0\%$, worst case $+133\%$;
        EnKF state$+$RTPS median $+68.7\%$. \good{ESMDA has the best
        typical regret and the hybrid the best worst case}, and the hybrid is
        the only variant that returns a usable estimate of \emph{both} halves
        in all three matched cases. If one configuration has to serve all
        regimes, that combination is the argument for it.
\end{description}

\subsubsection{Open issues and caveats}
\label{sec:cmp:caveats}

\begin{enumerate}[nosep,leftmargin=1.6em]
  \item \textbf{Unequal information.} The ESMDA rows work from $96$
        aggregated observations per window; the filter rows from $720$ raw
        ones. No run controls for this, so the ESMDA-versus-filter contrast
        confounds \emph{algorithm} with \emph{observation processing}. The
        turbulent-inflow result (ESMDA best on held-out sensors from less
        data) is robust to the confound in sign; the laminar result (filter
        best) is not.
  \item \textbf{The two observation-interval sweeps are different
        experiments} (caveat~3, \S\ref{sec:cmp:scope}), and they disagree:
        ESMDA's angle RMSE on \code{inflow\_turb} is flat across
        $7.5/15/30$~s ($2.42/2.31/2.35^\circ$) while the hybrid's is monotone
        ($1.94/2.61/2.86^\circ$). Both effects are $\leq 1^\circ$, an
        order of magnitude below the case-to-case differences, and the state
        metrics move by under $4\%$ in both campaigns. Neither sweep supports a
        recommendation.
  \item \textbf{One seed per configuration.} Every run uses \code{seed=42} and
        no configuration was repeated. Differences below roughly $10\%$ ---
        for example H13 $1.246$ versus F12 $1.261$ on
        \code{periodic} --- are not resolvable, and are described as ties
        above.
  \item \textbf{Six of $45$ matched cells are empty.} Both PALM-truth
        \code{inflow} cells failed in all three campaigns, at PALM truth
        generation rather than in assimilation, so the laminar model-error
        comparison does not exist. Two of the three columns of
        Table~\ref{tab:cmp:master-palm} are therefore missing a row that the
        matched table has.
  \item \textbf{Calibration is compared across instruments.} ESMDA
        contributes no $\chi^2$ and the filters contribute no $O_N$; the
        cross-method calibration statement rests on the $z$-scores alone.
        Every run also carries
        \code{caveat: no\_representativeness\_error}, so the absolute
        over-confidence is common-mode and only the \emph{ordering} between
        methods should be read from Figure~\ref{fig:cmp:calibration}.
  \item \textbf{Two RMSE definitions are in circulation.} The tables use the
        3-component vector RMSE; every figure produced by the run drivers
        (including the panels of Figure~\ref{fig:cmp:valid-ud-turb}) prints a
        magnitude-only $\Umag$ RMSE that is roughly half as large. They rank
        identically but must never be quoted side by side.
  \item \textbf{Wall times are not a clean cost model.} They mix truth models,
        cases and solver speeds, and all runs shared the same 10-process
        machine sequentially rather than being benchmarked in isolation.
  \item \textbf{The state-only rows are incomplete by design.} The filtering
        campaign ran state-only variants only for uDALES truth, and the
        uninflated variant only at \code{locnone}, so several cells of
        Table~\ref{tab:cmp:master-ud} compare four methods and others five.
\end{enumerate}

\subsection{At a glance}
\label{sec:cmp:glance}

\begin{table}[H]
  \centering
  \caption{Method $\times$ regime verdicts, matched model and localization on unless stated.
  \good{Green} = recommended for that regime, \bad{red} = avoid. Entries in
  black are usable but dominated by another method at equal or lower cost.
  Verdicts follow Tables~\ref{tab:cmp:master-ud}--\ref{tab:cmp:cost}.}
  \label{tab:cmp:glance}
  \footnotesize
  \begin{adjustbox}{max width=\textwidth}
  \begin{tabular}{@{}l llll@{}}
  \toprule
  method & \code{inflow} (laminar) & \code{inflow\_turb} & \code{periodic}
  & PALM truth (model error) \\
  \midrule
  ESMDA
    & E11: stalls, $0.748$ held-out
    & \good{E12: best, $0.365$, $2.31^\circ$}
    & E13: flat, $1.408$, $13.62^\circ$
    & \good{E2/E3: state OK, params wrong} \\
  EnKF state, no infl.
    & \bad{F14: $\chi^2 = 3.98$}
    & \bad{F17: $\chi^2 = 8.79$}
    & \bad{F20: $\chi^2 = 27.7$}
    & not run \\
  EnKF state$+$RTPS
    & F8: calibrated, no params
    & F10: dominated, $0.753$
    & \good{F12: ties best state, $1.261$}
    & not run \\
  EnKF joint$+$RTPS
    & \good{F13 (no loc): best, $0.319$, $3.90^\circ$}
    & F9: angle only, $3.36^\circ$
    & \bad{F11: degrades state \& params}
    & \bad{F3: diverges on periodic} \\
  Filter smoothing
    & H11: dominated, $0.593$
    & H12: close 2nd, $0.497$, $2.61^\circ$
    & \good{H13: best state $+$ safe params}
    & \good{H2/H3: most stable, $2.28\times$} \\
  \bottomrule
  \end{tabular}
  \end{adjustbox}
\end{table}

\noindent
Condensed to one sentence: \good{on this case the dominant error source
selects the method} --- boundary forcing $\rightarrow$ ESMDA, interior
phase error $\rightarrow$ joint EnKF, unidentifiable forcing $\rightarrow$
state-only EnKF or filter smoothing, model error $\rightarrow$ nothing
recovers the parameters and the state floor is method-independent --- and
\bad{no method in this set is simultaneously accurate and calibrated}.

\endgroup
